\documentclass[12pt,a4paper]{article}

\usepackage[utf8]{inputenc}
\usepackage[spanish]{babel}
\usepackage{amsmath, amssymb}
\usepackage{graphicx}
\usepackage{float}
\usepackage{booktabs}
\usepackage{geometry}
\geometry{margin=2.5cm}

\begin{document}
	
	% --- INICIO DE LA CARÁTULA ---
	\begin{titlepage}
		\centering
		{\bfseries\LARGE Universidad Nacional de Córdoba \par}
		{\scshape\Large Facultad de Ciencias Exactas, Físicas y Naturales \par}
		\vspace{2cm}
		%{\includegraphics[width=0.3\textwidth]{logo_unc}\par} % Si tenés el logo, si no comentá esta línea
		\vspace{1.5cm}
		{\bfseries\huge Sistemas de Computación \par}
		\vspace{1cm}
		{\Large Trabajo Practico 1\par}
		\vspace{2cm}
		
		{\Large \textbf{Integrantes:}} \par
		\vspace{0.5cm}
		{\large Mateo Bernardi - Ing. Comp.\par}
		{\large Gustavo Regnicoli - Ing Electr.\par}
			{\large Federico Schreiner - Ing Electr.\par}
		
		\vfill
		{\Large Año 2026 \par}
	\end{titlepage}
	% --- FIN DE LA CARÁTULA ---
	
	\newpage
	\tableofcontents % Esto genera el índice automáticamente
	\newpage
	
	\section{Introducción}
	El presente informe detalla el análisis de benchmarks de terceros para la toma de decisiones de hardware y la medición de performance de código propio mediante herramientas de profiling.
	
	\section{Parte 1: Benchmarks y Selección de Hardware}
	
	\subsection{Lista de Benchmarks Útiles}
	% Acá debés listar los benchmarks que investiguen (ej: Geekbench, Cinebench, Phoronix Test Suite)
	\begin{itemize}
		\item \textbf{Benchmark 1:} Descripción y utilidad para tareas de ingeniería.
		\item \textbf{Benchmark 2:} Aplicación en compilación y cálculo numérico.
	\end{itemize}
	
	\subsection{Relación Tarea-Benchmark}
	A continuación, se presenta la tabla de dos entradas que vincula las tareas diarias con el benchmark que mejor representa su carga de trabajo:
	
	\begin{table}[H]
		\centering
		\caption{Mapeo de Tareas Diarias a Benchmarks Específicos}
		\begin{tabular}{@{}ll@{}}
			\toprule
			\textbf{Tarea Diaria} & \textbf{Benchmark Representativo} \\ \midrule
			Compilación de código (C++/PlatformIO) & Timed Linux Kernel Compilation \\
			Simulación de Circuitos (SPICE) & SPECrate (Floating Point) \\
			Renderizado / Procesamiento de Imágenes & Cinebench R23 \\
			Uso general y navegación & PCMark 10 \\ \bottomrule
		\end{tabular}
	\end{table}
	
	\subsection{Análisis de Rendimiento: Compilación de Kernel Linux}
	Basado en los datos de \textit{OpenBenchmarking.org}, comparamos el tiempo de compilación del Kernel (menor es mejor):
	
	\begin{itemize}
		\item \textbf{Intel Core i5-13600K:} [Insertar valor aquí] segundos.
		\item \textbf{AMD Ryzen 9 5900X:} [Insertar valor aquí] segundos.
	\end{itemize}
	
	% Para la aceleración usá la fórmula: Speedup = T_referencia / T_nuevo
	\subsection{Cálculo de Aceleración (Speedup)}
	Utilizando el AMD Ryzen 9 7950X como base, la aceleración obtenida es:
	\begin{equation}
		S = \frac{T_{base}}{T_{7950X}} = \dots
	\end{equation}
	
	\section{Parte 2: Time Profiling}
	% Acá van tus capturas de pantalla
	\subsection{Realización del Tutorial}
	\begin{figure}[H]
		\centering
		% \includegraphics[width=0.8\textwidth]{captura_profiling.png}
		\caption{Captura de pantalla de la herramienta de profiling en ejecución.}
	\end{figure}
	
	\subsection{Conclusiones sobre el uso del tiempo}
	[Escribir aquí el análisis de qué funciones consumen más CPU].
	
	\section{Resultados}
	% Recordá que esto debe ir sincronizado pero con tus datos propios
	Mis mediciones individuales son: \dots
	
\end{document}